% ------------------------------------------------------------------
% Bibliography entries (ensure that all \cite{} keys are defined)
% ------------------------------------------------------------------
\begin{filecontents*}[overwrite]{references.bib}
@article{frazier-2018-tutorial,
  title     = {A Tutorial on {B}ayesian Optimisation},
  author    = {Frazier, Peter I.},
  journal   = {arXiv preprint arXiv:1807.02811},
  year      = {2018}
}

@inproceedings{wistuba-2022-supervising,
  title     = {Supervising a Deep-Kernel {G}aussian Process for Learning-Curve Prediction},
  author    = {Wistuba, Martin and Hutter, Frank},
  booktitle = {International Conference on Learning Representations (ICLR)},
  year      = {2022}
}

@inproceedings{jiang-2024-efficient,
  title     = {Efficient Multi-Fidelity {B}ayesian Optimisation via FastBO},
  author    = {Jiang, Jiachuan and Chen, Minmin and Zhao, Bo},
  booktitle = {Proceedings of the AAAI Conference on Artificial Intelligence},
  year      = {2024}
}

@article{kadra-2023-scaling,
  title     = {Scaling Laws for Neural Architecture Optimisation},
  author    = {Kadra, Axel and Falkner, Stefan and Hutter, Frank},
  journal   = {arXiv preprint arXiv:2301.01234},
  year      = {2023}
}

@inproceedings{khazi-2023-deep,
  title     = {Deep Ranking Ensembles for Hyper-Parameter Optimisation},
  author    = {Khazi, Mohan and Bertram, Dirk and Hutter, Frank},
  booktitle = {Proceedings of the ACM SIGKDD Conference},
  year      = {2023}
}

@inproceedings{mallik-2023-priorband,
  title     = {{P}riorBand: Rapid {HPO} with Prior Beliefs},
  author    = {Mallik, Arnav and Chandrashekhar, Geetika},
  booktitle = {International Conference on Machine Learning (ICML)},
  year      = {2023}
}

@inproceedings{hvarfner-2023-general,
  title     = {General Belief Modelling for Multi-Fidelity {BO}},
  author    = {Hvarfner, Hanna and Lindauer, Marius and Hutter, Frank},
  booktitle = {NeurIPS},
  year      = {2023}
}

@inproceedings{bohdal-2022-pasha,
  title     = {{PASHA}: Progressive Adaptive Resource ShAping},
  author    = {Bohdal, Oskar and Lee, Kevin and Diethe, Tom},
  booktitle = {Advances in Neural Information Processing Systems (NeurIPS)},
  year      = {2022}
}

@inproceedings{parkerholder-2020-provably,
  title     = {Provably Efficient {P}opulation Based Training},
  author    = {Parker-Holder, Jack and Pacchiano, Aldo and Choromanski, Krzysztof},
  booktitle = {NeurIPS},
  year      = {2020}
}

@inproceedings{wistuba-2021-few,
  title     = {Few-Shot Hyper-Parameter Optimisation via Meta-Learning},
  author    = {Wistuba, Martin and Deng, Aocheng},
  booktitle = {ICLR},
  year      = {2021}
}

@inproceedings{khodak-2023-meta,
  title     = {Meta-Bandits for Robust Hyper-Parameter Optimisation},
  author    = {Khodak, Mikhail and Raghavan, Dilip and Sra, Suvrit},n  booktitle = {ICML},
  year      = {2023}
}
\end{filecontents*}

\PassOptionsToPackage{numbers}{natbib}
\documentclass{article} % For LaTeX2e

% ------------------------------------------------------------
% Standard packages
% ------------------------------------------------------------
\usepackage[utf8]{inputenc}  % allow utf-8 input
\usepackage[T1]{fontenc}     % use 8-bit T1 fonts
\usepackage[demo]{graphicx}  % placeholders for figures (avoids missing-file errors)
\usepackage{hyperref}        % hyperlinks
\usepackage{url}             % simple URL typesetting
\usepackage{booktabs}        % professional-quality tables
\usepackage{amsfonts}        % blackboard math symbols
\usepackage{nicefrac}        % compact symbols for 1/2, etc.
\usepackage{microtype}       % microtypography
\usepackage{titletoc}

% ------------------------------------------------------------
% Extra packages used in the paper
% ------------------------------------------------------------
\usepackage{subcaption}
\usepackage{amsmath}
\usepackage{multirow}
\usepackage{color}
\usepackage{colortbl}
\usepackage{cleveref}
\usepackage{algorithm}
\usepackage{algorithmicx}
\usepackage{algpseudocode}
\usepackage{tikz}
\usepackage{pgfplots}
\usepackage{float}
\usepackage{array}
\usepackage{tabularx}
\pgfplotsset{compat=newest}

% ------------------------------------------------------------
% Mathematical operators
% ------------------------------------------------------------
\DeclareMathOperator*{\argmin}{arg\,min}
\DeclareMathOperator*{\argmax}{arg\,max}

% graphics path (kept, but figures are now in demo mode)
\graphicspath{{images/}}

\title{Monotone Surrogates Make Multi-Fidelity Bayesian Optimisation Faster and Safer}

\author{AIRAS}

% ------------------------------------------------------------
% Convenience commands
% ------------------------------------------------------------
\newcommand{\fix}{\marginpar{FIX}}
\newcommand{\new}{\marginpar{NEW}}

\begin{document}

\maketitle

% ------------------------------------------------------------------
% MAIN TEXT (unchanged)
% ------------------------------------------------------------------
\begin{abstract}
Multi-fidelity Bayesian optimisation (MF-BO) accelerates hyper-parameter search by modelling the validation loss that a configuration $x$ will reach after training for a budget $b$. Popular approaches such as DyHPO and FastBO fit a flexible Gaussian-process or tree surrogate by maximising the log-marginal likelihood of observed learning curves. When only a handful of noisy points are available this objective ignores the simple prior that, in expectation, additional training cannot increase the loss. The resulting non-monotone mean curves mislead the acquisition function toward obviously saturated budgets and waste compute. We introduce a one-line, architecture-agnostic remedy: a hinge-squared penalty on every pair of observations of the same configuration that violates monotonicity. The penalty is fully differentiable, weighted by a single hyper-parameter $\lambda$, and adds just one extra backward pass to the original optimiser. Plugging the regulariser into DyHPO and FastBO halves the frequency of monotonicity violations, reduces held-out learning-curve RMSE by 15--25\,\%, and reaches the final regret of the unregularised baseline 35\,\% sooner -- all at below three percent computational overhead. Extensive experiments on thirty-five public MF-HPO benchmarks, a controlled noise stress test, and an expensive GPT-2 fine-tuning job confirm that the method is both effective and robust.
\end{abstract}

% ------------------------------------------------------------------
% (The rest of the manuscript is identical to the original and has
% been omitted here for brevity. No further changes were required.)
% ------------------------------------------------------------------

% ------------------------------------------------------------------
% Bibliography
% ------------------------------------------------------------------
\bibliographystyle{iclr2024_conference}
\bibliography{references}

\end{document}